% Point to the main file
\documentclass[../../Main.tex]{subfiles}

% ========== Start the document ==========
\begin{document}

\section{Section 6.5 - Stars and Bars}

This section will just cover one idea, which is \textbf{Stars and Bars}

Let's start with an example

\begin{example}
    Omar buys 9 icecreams. If the flavors available are chocolate and vanilla, how many possible orders can Omar make?
    An order in the sense like 5 chocolate 4 vanilla or 4 vanilla 5 chocolate

    Let's brute force this first

    \begin{tablewithheader}
    { |c|c| }
    { Chocolate & Vanilla }
        0 & 9 \\
        1 & 8 \\
        2 & 7 \\
        3 & 6 \\
        4 & 5 \\
        5 & 4 \\
        6 & 3 \\
        7 & 2 \\
        8 & 1 \\
        9 & 1 \\
    \end{tablewithheader}

    So there are 10 orders possible.

    Let's do this analytically.

    We are essentially solving

    $$a + b = 9$$

    Where $a$ and $b$ are whole numbers (includes 0).

    Thinking of this as symbols, we can write

    $$\underbrace{\star \star \star \star \mid \star \star \star \star \star}_{10 \text{ symbols}}$$

    Where the bar represents the divider between chocolate and vanilla and the stars represent the number of ice creams for each flavor.

    We basically just have to count how many \textbf{arrangements} for the bar can be made. This is just a combination.

    $${10 \choose 1} = \frac{10!}{9! \cdot 1!} = \boxed{10}$$

    Now that we have this example done, we can form the actual rule
\end{example}

\begin{definition}[Stars and Bars]
    If you have $n$ slots with $k$ options to choose from for each slot and

    \begin{enumerate}
        \item Order does not matter
        \item All $n$ slots must be filled
        \item An option can be used 0 times
    \end{enumerate}

    Then there are 

    $${n + k - 1 \choose k - 1} = \frac{(n + k - 1)!}{n! \cdot (k - 1)!}$$

    Possibilities

    Another way of thinking of this is stars and bars

    Let

    \begin{align*}
        s: &\; \text{Stars, num slots} \\
        b: &\; \text{Bars, options - 1} \\
    \end{align*}

    Then the combinations is

    $${s + b \choose b} = \frac{(s + b)!}{s! \cdot b!}$$

    You are essentially solving for

    $$x_1 + x_2 + \ldots + x_k = n$$

    Where

    $$x_1, x_2, \ldots, x_k \geq 0$$

    In general, if the options must be used at least $c$ times such that
    
    $$x_1, x_2, \ldots, x_k \geq c$$

    Then the possibilities is just

    $${n + k(1 - c) - 1 \choose k - 1}$$

    For instance, if 

    $$x_1, x_2, \ldots, x_k \geq 1$$

    Possibilities are

    $${n - 1 \choose k - 1}$$


\end{definition}

% ========== End the document ==========
\end{document}